\section{Introduction}
\label{sec:introduction}
Speculative decoding accelerates autoregressive inference by proposing tokens with a draft and evaluating several positions with the target~\cite{leviathan2023}. Its appeal on consumer systems is practical: additional intra-request work can amortize target execution when memory and traffic do not permit large inter-request batches. Yet replacing scalar continuation with speculative execution changes numerical call shapes, retained state, and bookkeeping. A mathematically valid sampling rule does not by itself certify that a particular quantized implementation reproduces a selected greedy output or lowers the cost of doing so.

We consider systems in which an ordinary greedy execution defines the required reference. Examples include regression testing and repeatable application behavior under a registered runtime configuration. The requirement is token or endpoint equality to that execution, rather than a claim that its answer is functionally correct or that every floating-point value is identical. This makes two engineering decisions inseparable. The designer must determine where execution consistency needs to be enforced, then measure whether the qualified transformation is useful. Constraining the wrong boundary can repair one input and introduce another deviation. Constraining all computation can recover output while losing the intended acceleration.

The boundary is not determined by the visible token prefix alone. Two executions can have emitted the same tokens while retaining different numerical key/value (KV) contents or pending state. Their next decision can depend on that history, on the current execution recipe, or on an interaction between them. Conversely, changing a score need not change the immediate greedy decision. An unchanged decision at one verification block can still precede different downstream output. These distinctions determine whether an implementation needs to constrain state formation, current continuation, or both.

The cost comparison has a related ambiguity. A repair can improve the target path before any drafting occurs. Comparing only repaired speculation with the ordinary reference assigns that path benefit to the entire deployment change. It does not isolate the value of speculation. A matched target-only path, using the repaired target execution without a draft, is therefore a second paid comparator. Neither comparator can be replaced by accepted-token fraction, target-call count, or isolated draft throughput.

Our cases show why both decisions matter. An RTX continuation intervention repairs one task and regresses another while preserving aggregate reference agreement. A later phase-scoped attention implementation qualifies full responses on 24 inputs unused in the accessible project history. Its resident-request speed is $1.742\times$ ordinary decoding, but matched draft-free execution already gives $1.683\times$; speculation adds $1.035\times$ on that path. A separate B570 implementation preserves the registered endpoint yet is substantially slower than both paid draft-free comparators. These outcomes support a system design study built around three questions:
\begin{enumerate}
\item \RQtag{1}, diagnosis: does a deviation follow retained numerical state, the current recipe, or their interaction?
\item \RQtag{2}, repair: within a defined implementation family, which phase constraints are sufficient to preserve the reference?
\item \RQtag{3}, cost: when qualification holds, what benefit remains against ordinary decoding and against a matched target path?
\end{enumerate}

The contributions are concrete. First, controlled natural-state crosses with own-path shams expose different local dependencies on CUDA and Vulkan; a separate same-state B570 fork study shows why immediate agreement cannot replace trajectory qualification. Second, an RTX query-layout transformation preserves batched projection and feed-forward work, with a four-switch factorial locating a sufficient target-phase scope in the tested family. Third, ordinary/matched/candidate costs and qualified counterexamples change the engineering decision about adding speculation, increasing draft width, retaining a guard, or optimizing only the draft.

These contributions concern an implemented runtime and finite qualification, building on established numerical-consistency and verification principles. Diagnosis remains an offline experimental procedure; the model-family and measurement limits are stated in Section~\ref{sec:implications}.

Sections~\ref{sec:runtime} and~\ref{sec:methods} define the runtime and contracts. Sections~\ref{sec:diagnosis}--\ref{sec:cost} answer the three questions. Section~\ref{sec:implications} states the resulting design decisions and limits. The independent supplement gives campaign details, historical results, failures and missingness; the author source package supplies arithmetic and figure reconstruction.
